% Generated draft: requires mathematical and source review.
\documentclass[11pt]{article}
\usepackage[margin=25mm]{geometry}
\usepackage{fontspec}
\setmainfont{Latin Modern Roman}
\usepackage{amsmath,amssymb,mathtools,unicode-math}
\setmathfont{Latin Modern Math}
\usepackage{enumitem,needspace,xurl,hyperref}
\setlength{\parindent}{0pt}
\setlength{\parskip}{5pt}
\setlength{\emergencystretch}{3em}
\newcommand{\N}{\mathbb{N}}
\newcommand{\Z}{\mathbb{Z}}
\newcommand{\Q}{\mathbb{Q}}
\newcommand{\R}{\mathbb{R}}
\newcommand{\C}{\mathbb{C}}
\newcommand{\F}{\mathbb{F}}
\newcommand{\A}{\mathbb{A}}
\newcommand{\PP}{\mathbb{P}}
\newcommand{\E}{\mathbb{E}}
\newcommand{\eps}{\varepsilon}
\newcommand{\dd}{\,\mathrm{d}}
\begin{document}
\section*{UnsolvedMath ID: 30004732}
\textbf{Source problem code:} catalogue entry
\subsection{Definitions and mathematical statement}
A Hadamard manifold is a complete, simply connected Riemannian manifold of nonpositive sectional curvature. Let $b\in[1,\infty)$, let $X$ be such a manifold satisfying
\[
-b^2\leq K_X(\sigma)\leq-1
\]
for every tangent two-plane $\sigma$, and let $\Gamma\leq\operatorname{Isom}(X)$ be a finitely generated discrete subgroup. For a base point $o\in X$, define its critical exponent by
\[
\delta_\Gamma=\inf\left\{s\geq0:\sum_{\gamma\in\Gamma}e^{-s d_X(o,\gamma o)}<\infty\right\}.
\]
This definition is independent of $o$. A torsion element is an element of finite order.

For a precise group-theoretic reading of the cusp assertion, let $\partial_\infty X$ denote the visual boundary, and let $\mathcal P(\Gamma)$ be the set of points of $\partial_\infty X$ fixed by parabolic elements of $\Gamma$. Here a parabolic element has no fixed point in $X$ and exactly one fixed point in $\partial_\infty X$. Interpret finiteness of cusps as
\[
\#\bigl(\Gamma\backslash\mathcal P(\Gamma)\bigr)<\infty.
\]
The supplied record does not define ``cusp''; whether it intends precisely these parabolic classes or a more restrictive geometric notion of cusp end requires review.

\textbf{Supplied conjectural assertion.} For every $b,X,\Gamma$ as above, does $\delta_\Gamma\leq2$ imply both finiteness of cusps and
\[
\exists M\in\N\ \forall\gamma\in\Gamma:\quad
\operatorname{ord}(\gamma)<\infty\ \Longrightarrow\ \operatorname{ord}(\gamma)\leq M?
\]
The bound $M$ may depend on the given action; no universal bound across all such groups is asserted.

\textbf{Additional existence assertion.} Do there exist sequences $(b_k,X_k,\Gamma_k)_{k\geq1}$ such that each $X_k$ is a four-dimensional Hadamard manifold, each $\Gamma_k\leq\operatorname{Isom}(X_k)$ is finitely generated and discrete, and
\[
1\leq b_k<\infty,\qquad
-b_k^2\leq K_{X_k}\leq-1,\qquad
\delta_{\Gamma_k}>2,\qquad
\lim_{k\to\infty}\delta_{\Gamma_k}=2,\qquad
\lim_{k\to\infty}b_k=1?
\]
Here ``from above'' is expressed by $\delta_{\Gamma_k}>2$ and convergence, without imposing monotonicity. The word ``such'' in the original sequence assertion cannot also retain $\delta_{\Gamma_k}\leq2$. The display retains the geometric and group hypotheses while using the explicitly requested approach from above. The record does not specify whether these examples must additionally have infinitely many cusps or unbounded torsion orders; no such additional requirement is inserted.
\subsection{Short English statement}
The supplied conjecture asks whether critical exponent at most $2$ forces finitely many cusps and a bound on torsion orders for every finitely generated discrete action on a negatively pinched Hadamard manifold. It also asks for four-dimensional examples whose critical exponents approach $2$ from above while their curvature pinching constants approach $1$. The meaning of cusp and the intended additional properties of the sequence require clarification from the source.
\subsection{Sources}
\begin{itemize}
\item Supplied local UnsolvedMath registry record, permanent ID $30004732$, exact source problem code \texttt{catalogue entry}, title ``Critical-Exponent Geometry of Negatively Pinched Hadamard Quotients.'' The record belongs to the Oberwolfach Reports Open Problems collection, gives proposed year $2021$, and names no proposer.
\item Supplied report citation: \emph{Differentialgeometrie im Grossen} (2022), Oberwolfach Reports, DOI 10.4171/owr/2021/32, pp.~1685--1734. \url{https://doi.org/10.4171/owr/2021/32}. The background identifies the extracted problem as Conjecture~6. The supplied record labels its statement \texttt{corrected\_verified} and reports checking it against the official report; this draft does not independently verify that claim.
\item Supplied literature reference labeled ``orbifold geometrization'': \url{https://annals.math.princeton.edu/2005/162-1/p04}.
\item Supplied literature reference labeled ``small critical exponent subcase'': \url{https://arxiv.org/abs/2008.12365}.
\item Literature assessment dated exactly \texttt{2026-08-22}: the supplied curation labels the problem partially solved, states that orbifold geometrization is proved, and attributes to Liu--Wang the cusp-finiteness subcase $\delta_\Gamma<1$. It states that the extracted Conjecture~6 for $\delta_\Gamma\leq2$, the full bounded-torsion assertion, and the four-dimensional sequence assertion remain unresolved. These are supplied dated curation claims, not independently verified mathematical or literature assessments. The background explicitly requests independent verification of time-sensitive status and cited claims.
\end{itemize}
\end{document}
